% BEGIN VOL08-EXPANSION VIII24
\section{Supplementary worked examples}
\label{sec:viii24-pedagogy}

\begin{example}[Cone of the identity is acyclic]\label{ex:viii24-ped-01}
For a chain map $f:C_*\to D_*$, the mapping cone has groups
\[
\operatorname{Cone}(f)_n=D_n\oplus C_{n-1}.
\]
When $f=\operatorname{id}_C$, the cone is chain contractible, reflecting that the topological cone on an identity attachment has no residual relative homology.
\end{example}

\begin{example}[Cone of the zero map]\label{ex:viii24-ped-02}
For $f=0:C\to D$, the differential does not couple the two summands, so
\[
\operatorname{Cone}(0)\cong D\oplus \Sigma C
\]
as chain complexes. Consequently $H_n(\operatorname{Cone}(0))\cong H_n(D)\oplus H_{n-1}(C)$.
\end{example}

\begin{example}[Quasi-isomorphism criterion]\label{ex:viii24-ped-03}
The cone fits into a short exact sequence
\[
0\to D\to\operatorname{Cone}(f)\to\Sigma C\to0.
\]
Its long exact homology sequence shows that $f$ is a quasi-isomorphism exactly when $\operatorname{Cone}(f)$ is acyclic.
\end{example}

\section{Graded supplementary exercises}
The following exercises progress from direct computation and construction to proof, hypothesis testing, application, and synthesis.

\subsection*{Standard computations and constructions}

\begin{exercise}[Cone group]\label{exr:viii24-ped-01}
Write $\operatorname{Cone}(f)_n$.
\end{exercise}
\begin{hint}
Shift the source down by one degree.
\end{hint}
\begin{solution}
$\operatorname{Cone}(f)_n=D_n\oplus C_{n-1}$.
\end{solution}

\begin{exercise}[Cone differential]\label{exr:viii24-ped-02}
With homological grading, write a standard cone differential.
\end{exercise}
\begin{hint}
Include $f(c)$ in the $D$ component and a minus sign on $d_C$.
\end{hint}
\begin{solution}
$d(d,c)=(d_Dd+f(c),-d_Cc)$.
\end{solution}

\begin{exercise}[Zero map cone]\label{exr:viii24-ped-03}
Identify $\operatorname{Cone}(0)$.
\end{exercise}
\begin{hint}
The coupling term disappears.
\end{hint}
\begin{solution}
$\operatorname{Cone}(0)\cong D\oplus\Sigma C$.
\end{solution}

\begin{exercise}[Identity cone]\label{exr:viii24-ped-04}
What is the homology of $\operatorname{Cone}(\operatorname{id}_C)$?
\end{exercise}
\begin{hint}
The identity is a quasi-isomorphism.
\end{hint}
\begin{solution}
It is zero in every degree; the cone is in fact contractible.
\end{solution}

\begin{exercise}[Shifted homology]\label{exr:viii24-ped-05}
What is $H_n(\Sigma C)$ in terms of $H_*(C)$?
\end{exercise}
\begin{hint}
The shift moves degree by one.
\end{hint}
\begin{solution}
$H_n(\Sigma C)\cong H_{n-1}(C)$.
\end{solution}

\subsection*{Proofs}

\begin{exercise}[Check $d^2=0$]\label{exr:viii24-ped-06}
Verify the cone differential squares to zero.
\end{exercise}
\begin{hint}
Use $d_Df=fd_C$.
\end{hint}
\begin{solution}
Applying $d$ twice cancels the two $f(d_Cc)$ terms and leaves $d_D^2d$ and $d_C^2c$, both zero.
\end{solution}

\begin{exercise}[Short exact cone sequence]\label{exr:viii24-ped-07}
Prove $0\to D\to\operatorname{Cone}(f)\to\Sigma C\to0$ is a short exact sequence of complexes.
\end{exercise}
\begin{hint}
Use inclusion into the first summand and projection to the second.
\end{hint}
\begin{solution}
The maps are degreewise exact, and the cone differential is compatible with inclusion and projection under the standard shifted differential.
\end{solution}

\begin{exercise}[Quasi-isomorphism criterion]\label{exr:viii24-ped-08}
Prove $f$ is a quasi-isomorphism iff $\operatorname{Cone}(f)$ is acyclic.
\end{exercise}
\begin{hint}
Use the long exact homology sequence of the cone short exact sequence.
\end{hint}
\begin{solution}
Exactness identifies the connecting morphism with $f_*$ up to sign; all cone homology vanishes exactly when every $f_*$ is an isomorphism.
\end{solution}

\begin{exercise}[Homotopic maps and cones]\label{exr:viii24-ped-09}
Explain why chain-homotopic maps have chain-homotopy-equivalent cones.
\end{exercise}
\begin{hint}
Use the homotopy operator to shear one cone summand.
\end{hint}
\begin{solution}
A chain homotopy gives an explicit triangular chain isomorphism/homotopy equivalence between the two cone constructions, so their homology agrees.
\end{solution}

\subsection*{Counterexamples and hypothesis tests}

\begin{exercise}[Cone acyclic means map zero]\label{exr:viii24-ped-10}
Test: if $\operatorname{Cone}(f)$ is acyclic, then $f$ is the zero map.
\end{exercise}
\begin{hint}
Acyclic cone characterizes quasi-isomorphisms.
\end{hint}
\begin{solution}
False. The identity map has an acyclic cone and is usually far from zero.
\end{solution}

\begin{exercise}[Cone of zero always acyclic]\label{exr:viii24-ped-11}
Test: $\operatorname{Cone}(0)$ is always acyclic.
\end{exercise}
\begin{hint}
It contains $D$ and a shifted copy of $C$.
\end{hint}
\begin{solution}
False; its homology is $H_n(D)\oplus H_{n-1}(C)$.
\end{solution}

\begin{exercise}[Cone forgets f]\label{exr:viii24-ped-12}
Test: the cone differential depends only on $C$ and $D$, not on $f$.
\end{exercise}
\begin{hint}
Inspect the coupling term.
\end{hint}
\begin{solution}
False. The term $f(c)$ is precisely where the map is encoded.
\end{solution}

\subsection*{Applications and investigations}

\begin{exercise}[Relative model]\label{exr:viii24-ped-13}
Why do mapping cones model relative phenomena?
\end{exercise}
\begin{hint}
A cone kills the source by attaching a shifted copy.
\end{hint}
\begin{solution}
The cone packages a map together with a formal null-homotopy of its source, paralleling the topological mapping cone whose reduced homology enters relative exact sequences.
\end{solution}

\begin{exercise}[Testing quasi-isomorphisms]\label{exr:viii24-ped-14}
How can the cone reduce a many-degree isomorphism test to one homology computation?
\end{exercise}
\begin{hint}
Compute whether the cone is acyclic.
\end{hint}
\begin{solution}
Instead of checking each induced map directly, build the cone and verify that all its homology groups vanish.
\end{solution}

\subsection*{Challenge problems}

\begin{exercise}[Cone of multiplication]\label{exr:viii24-ped-15}
For $f:\mathbb Z\to\mathbb Z$ given by multiplication by $m$ between complexes concentrated in degree $0$, compute cone homology.
\end{exercise}
\begin{hint}
The cone is the two-term complex $\mathbb Z\xrightarrow{m}\mathbb Z$.
\end{hint}
\begin{solution}
For $m\ne0$, $H_0\cong\mathbb Z/m$ and $H_1=0$; for $m=0$, both expected shifted/source terms appear.
\end{solution}

\begin{exercise}[Cone synthesis]\label{exr:viii24-ped-16}
Describe how mapping cones unify cokernels, shifts, exact sequences, and quasi-isomorphisms.
\end{exercise}
\begin{hint}
Use the cone short exact sequence.
\end{hint}
\begin{solution}
The cone joins $D$ to a shifted $C$, yields a long exact sequence whose connecting map is $f_*$, reduces to a direct sum for $f=0$, and is acyclic precisely for quasi-isomorphisms.
\end{solution}

% END VOL08-EXPANSION VIII24
